\newpage\section{Conventions and the common canonical element}

Put \(r=m,\ N=2m+1,\ I=\{s_1,\ldots,s_{N-1}\}\), \(w_I=w_0\), and
\(V=I\setminus\{s_m,s_{m+1}\}\). Both papers compose periodic permutations by
\((xy)(i)=x(y(i))\). Right multiplication by \(s_i\) swaps positions \(i,i+1\).
For \(s_0\) the new first and last entries are \(w(N)-N,w(1)+N\).
The increasing window selects the shortest representative in \(W/W_I\);
the increasing inverse window selects the shortest representative in
\(W_I\backslash W\). We work in the extended group. Our construction has degree zero.

\paragraph{Parameters.}
Both use \(H_s^2=1+(v^{-1}-v)H_s,\ C_s=H_s+v\). With
\(T_w=v^{-\ell(w)}H_w\), the coefficient of \(H_y\) in \(C_w\) is
\(h_{y,w}(v)=v^{\ell(w)-\ell(y)}P_{y,w}(v^{-2})\).
We use \(q=v^{-2}\) in our flag counts. Bowman, de la Fuente, Libedinsky and Plaza use \(q=v^2\) for their Gaussian polynomials (B, \S3.1). Gaussian reciprocity contributes the factor
\(v^{-2b(a-b)}\) when changing \(\gb{a}{b}_{v^{-2}}\) to
\(\gb{a}{b}_{v^2}\). For \(L=\ell(w_P)\), we write \(C_P,h_P\) for their
\(C_{w_P},\pi_P\), respectively:
\[
 C_P=\sum_{z\in W_P}v^{L-\ell(z)}H_z,\qquad
 h_P=\sum_{z\in W_P}v^{L-2\ell(z)},\qquad C_P^2=h_PC_P .
\]

\paragraph{Sides and bases.}
Our web maps the right ideal \(C_I\mathcal H\) to \(C_V\mathcal H\).
Evaluate on \(C_I\), then apply the linear anti-involution
\(H_w^*=H_{w^{-1}}\), which fixes \(v\) and reverses products. The result
lies in their right ideal \(C_I\mathcal H\), with right \(V\)-spherical character.
For a double coset \(D=W_VdW_I\) with longest element \(z\), its standard
map goes to \(v^{\ell(z)}\sum_{w\in D^{-1}}T_w\), and its canonical map
to \(C_{z^{-1}}\). Their bases are \(M_b=C_IH_b\) and
\(\mathbf C_b=C_{w_Ib}\), for \(b\in{}^IW\), and
\([M_a]\mathbf C_b=h_{w_Ia,w_Ib}\).
At our splitter the inverted standard is \(M_e=C_I\).
Thus for \emph{all} \(j\), with \(b=u_{j;m}^{-1}\),
\[
 \big[[B_m]_v\big]\theta_{A_{j;m}}=[M_e]\mathbf C_b
   =v^{D_{j;m}}P_{w_0,C_{j;m}}(v^{-2}),\qquad
 D_{j;m}=2j(N-j)-2(m-j)-3 .
\]
Here \(C_{j;m}=u_{j;m}w_0\) is a permutation, while \(C_w\) denotes a
basis vector. See A Definitions 2D.1--3, 2E.1--2, Lemma 2E.7,
Proposition 2E.9 and \S3B; B \S2.1, Lemma 2.1.

\paragraph{The diagonal.}
Write \(\nu=(1^m,0,(-1)^m)\), \(\sigma=(m,m+2)\), and
\(t_\nu(i)=i+N\nu_i\) on \(1\le i\le N\). In our increasing window at \(j=m\)
the ordered residues are
\(1,\ldots,m-1,m+2;\ m+1;\ m,m+3,\ldots,N\) with shifts \(-1,0,1\)
on the respective groups. Therefore \(u_{m;m}=\sigma t_{-\nu}\).
Its inverse \(t_\nu\sigma\) has window
\([N+1,\ldots,N+m-1,1-m,m+1,N+m,2-m,\ldots,0]\), the window of their element \(\beta_m\).
As \(w_I\beta_m=C_{m;m}^{-1}\), inversion identifies the ordinary KL polynomials
and the full Hecke vectors satisfy
\((\theta_{A_{m;m}}(C_I))^*=\mathbf C_{\beta_m}\).
The dot action gives \(\beta_m\cdot_p0=p\nu+\sigma\rho-\rho
=p\nu-2(\alpha_m+\alpha_{m+1})\), hence
\[
 \lambda_{m;m}(p)=2\varpi_{m-1}+
        (p-2)(\varpi_m+\varpi_{m+1})+2\varpi_{m+2}.
\]
They include the factor \(p\) in the translation label. See A Lemma 6A.2
and B Lemma 2.16.
